\begin{proofbox}
	\textbf{\textcolor{CProof}{思路提示}}

	利用抛物线定义将焦半径转化为点到准线的距离，再通过坐标差直接计算，避免两点间距离公式。

	\medskip
	\textbf{\textcolor{CProof}{正式推导}}
	\begin{description}[style=nextline, leftmargin=2.8em, labelsep=0.5em, itemsep=0.55em]
		\item[\textbf{步骤一（明确几何要素）：}]
			由抛物线标准方程 $y^2=2px$（$p>0$），得焦点 $F$ 和准线 $l$：
			\[
				F\left(\frac{p}{2},0\right),\quad l:x=-\frac{p}{2}.
			\]
			\par\textbf{理由：}
			标准方程的几何参数直接对应。

		\item[\textbf{步骤二（定义转化）：}]
			根据抛物线定义，点 $P(x_0,y_0)$ 在抛物线上，则到焦点的距离等于到准线的距离：
			\[
				|PF| = d(P,l).
			\]
			\par\textbf{理由：}
			抛物线定义：平面内到定点与定直线距离相等的点的轨迹。

		\item[\textbf{步骤三（代数化简）：}]
			点到直线 $x=-\frac{p}{2}$ 的距离为横坐标差，即
			\[
				d(P,l) = \left|x_0 + \frac{p}{2}\right|.
			\]
			因点 $P$ 在抛物线上，$x_0 \ge 0$，所以
			\[
				|PF| = x_0 + \frac{p}{2}.
			\]
			\par\textbf{理由：}
			抛物线 $y^2=2px$ 开口向右，抛物线上点横坐标非负，故绝对值可去。
	\end{description}

	\medskip
	\textbf{\textcolor{CProof}{结论回扣}}

	\textbf{该证明通过抛物线定义将距离转化，避免了两点间距离公式的复杂运算，说明焦半径公式在求解抛物线问题时的简洁性。}
\end{proofbox}
